\section{La bifurcación de rutas entre el proveedor de OpenAI y los proveedores que no son de OpenAI}

\subsection{Dos rutas de código distintas}

En el código fuente de Codex, según el tipo de proveedor, se sigue una ruta de Compact distinta:

\begin{itemize}
    \item \textbf{Proveedor de OpenAI}: invoca la API dedicada \texttt{response.compact}, una interfaz de compresión especializada a la que basta con enviar el historial de la conversación para obtener el contenido comprimido. El System Prompt y el Handoff Prompt se mantienen en el \textbf{extremo remoto} (del lado del servidor de OpenAI).
    \item \textbf{Proveedores que no son de OpenAI}: invocan la interfaz común \texttt{response.create}, e inyectan localmente el Compact Prompt y el Handoff Prompt.
\end{itemize}

\begin{knowledgebox}{La API response.compact}
Es la interfaz de compresión especializada que ofrece OpenAI. Quien la invoca solo necesita enviar el historial de la conversación (messages), y la API devuelve el resumen comprimido. No hace falta mantener los prompts de compresión del lado del cliente.
\end{knowledgebox}

\subsection{Resumen de la sección}

El proveedor determina dónde se mantiene el prompt de compresión: con OpenAI se sigue la ruta remota de \texttt{response.compact}, mientras que con el resto se sigue la ruta local de \texttt{response.create} con inyección explícita del prompt.

